\sectionTitle{Experience}{\faSuitcase}
\begin{experiences}
	\experience
		{Now} {Software Engineer}{ION GROUP}{}
		{March 2023} {
						\begin{itemize}
						\item Technical referent for two production products spanning multiple components, with backend systems managing databases of millions of rows and processing thousands of events daily.
						\item Sustained a near-zero incident rate on owned systems; acted as fastest responder when issues originated in upstream systems, given deeper system knowledge.
						\item Owned and maintained a critical Quarkus backend component containing core business logic.
						\item Led end-to-end feature development with increasing scope and business impact.
						\item Contributed to the early-stage development of a new product, from design to implementation.
						\item Collaborated with Product Management to translate customer needs into technical requirements.
						\item Contributed to architectural decisions and cross-team alignment.
						\end{itemize}
		}
		{Java, Quarkus, Jira, Confluence, git}
	\emptySeparator
	\experience
		{March 2023} {Data Engineer}{REPLY S.P.A.}{}
		{June 2022} {
						\begin{itemize}
						\item Built and deployed to production a daily pipeline ingesting Excel files and transforming them into Parquet tables following a bronze/silver/gold (medallion) architecture.
						\item Designed and implemented ETL pipelines ingesting data from external APIs into Azure Data Lake.
						\item Processed and refined raw data using Databricks and PySpark across multiple lake layers.
						\item Built Power BI dashboards consuming refined data from the lake for business reporting.
						\item Contributed to architectural decisions for a customer building their first Data Platform.
						\end{itemize}
		}
		{Microsoft Azure, PySpark, Python, Databricks, Power BI, git}
	\emptySeparator
	\experience
		{December 2021} {Thesis Researcher (Information Retrieval)}{UNIVERSITY OF GLASGOW}{}
		{October 2021} {
						\begin{itemize}
						\item Conducted a deep literature review on static pruning techniques applied to inverted indexes.
						\item Studied classical and neural information retrieval models, including BERT-based dense retrievers (ColBERT).
						\item Designed and experimented with static pruning strategies adapted to dense retrieval indexes.
						\item Published results at DocEng'23: 'Static Pruning for Multi-Representation Dense Retrieval'.
						\end{itemize}
						\href{https://dl.acm.org/doi/10.1145/3573128.3604896}{\textit{Static Pruning for Multi-Representation Dense Retrieval}} --- DocEng'23.
		}
		{PyTerrier, Jupyter Notebook, BERT, ColBERT}
	\emptySeparator
	\experience
		{June 2019} {Freelance Front-end Developer}{Graldev S.a.s.}{}
		{July 2018} {
						\begin{itemize}
						\item Developed and maintained UI screens for an iOS application using Swift and xCode.
						\item Integrated REST APIs to fetch and display server-side data within the app.
						\item Handled first interactions with stakeholders to gather requirements and report progress.
						\item Provided support to the production application, diagnosing and fixing reported issues.
						\end{itemize}
		}
		{Swift, xCode, git}
\end{experiences}
